\documentclass[11pt,a4paper]{article}

% ===== XeLaTeX + Korean =====
\usepackage{kotex}
\usepackage{fontspec}
\setmainfont{Times New Roman}
\setmainhangulfont{AppleMyungjo}
\setsanshangulfont{Apple SD Gothic Neo}

% ===== Layout =====
\usepackage[a4paper,top=18mm,bottom=18mm,left=20mm,right=20mm]{geometry}
\usepackage{fancyhdr}
\setlength{\headheight}{24pt}
\usepackage{enumitem}
\usepackage{mdframed}
\usepackage{array}

\setlength{\parindent}{1em}
\linespread{1.18}

% ===== Header / Footer =====
\pagestyle{fancy}
\fancyhf{}
\renewcommand{\headrulewidth}{0.6pt}
\fancyhead[L]{\sffamily\bfseries 2026 고1 영어 변형 — 지문 18}
\fancyhead[R]{\sffamily 도서 기부 행사 안내문}
\renewcommand{\footrulewidth}{0pt}
\fancyfoot[C]{\framebox[16mm][c]{\thepage}}

% ===== helpers =====
\newcommand{\qno}[1]{\par\vspace{8pt}\noindent{\sffamily\bfseries #1.}\ }
\newcommand{\instr}[1]{{\sffamily #1}\par\vspace{3pt}}
\newlist{choices}{enumerate}{1}
\setlist[choices]{label=,leftmargin=1.5em,itemsep=2pt,topsep=3pt,parsep=0pt}

\newmdenv[linewidth=0.6pt,roundcorner=3pt,innertopmargin=7pt,innerbottommargin=7pt,
          innerleftmargin=9pt,innerrightmargin=9pt,skipabove=5pt,skipbelow=5pt]{pbox}
\newcommand{\nemo}[1]{\fbox{\,#1\,}}
\newcommand{\ansline}{\par\vspace{2pt}\noindent\rule{\linewidth}{0.4pt}\par\vspace{8pt}\noindent\rule{\linewidth}{0.4pt}}

\begin{document}

\begin{center}
{\fontsize{15}{17}\selectfont\sffamily\bfseries [지문 18] 다음 글을 읽고 물음에 답하시오.}
\end{center}
\vspace{4pt}

% ===== 공통 지문 (원본 단어 유지 + 어법 네모) =====
\begin{pbox}
Dear Residents,\par
\hspace{1em}I am Trixie Mitchell, the director of the Riverside Community Center. I am pleased (A)\,\nemo{to invite / inviting} you to a book donation drive. We are collecting used books (B)\,\nemo{to support / supported} our reading programs for children and teenagers. By donating books, you can help young readers discover the joy of reading. Any books in good condition (C)\,\nemo{is / are} gladly accepted, including novels, poetry, and non-fiction books. You can bring your books to our community center during our operating hours. We would be grateful if you could join this meaningful event.\par
Sincerely,\par
Trixie Mitchell
\end{pbox}

% ===================== Q1 =====================
\qno{1}\instr{윗글의 목적으로 가장 적절한 것은?}
\begin{choices}
\item ① to announce the extended operating hours of the community center
\item ② to ask residents to donate used books for reading programs
\item ③ to recruit volunteers who will teach children how to read
\item ④ to promote a charity sale of secondhand books at low prices
\item ⑤ to invite residents to the grand opening of a new public library
\end{choices}

% ===================== Q2 =====================
\qno{2}\instr{윗글의 (A), (B), (C)의 각 네모 안에서 어법에 맞는 표현으로 가장 적절한 것은?}
\begin{center}
\renewcommand{\arraystretch}{1.25}
\begin{tabular}{c c c c}
 & (A) & (B) & (C) \\
① & to invite & to support & is \\
② & to invite & to support & are \\
③ & to invite & supported & are \\
④ & inviting & to support & is \\
⑤ & inviting & supported & are \\
\end{tabular}
\end{center}

% ===================== Q3 =====================
\qno{3}\instr{윗글의 내용과 일치하지 \underline{않는} 것은?}
\begin{choices}
\item ① The writer is responsible for running the Riverside Community Center.
\item ② The donation campaign aims to benefit young readers' reading habits.
\item ③ Worn-out or damaged books are unlikely to be welcomed.
\item ④ Factual books such as biographies fall outside what is being collected.
\item ⑤ Donors are expected to drop off their books when the center is open.
\end{choices}

% ===================== Q4 (서술형) =====================
\qno{4}\instr{[서술형] 윗글을 읽고 다음 물음에 답하시오.}
\vspace{2pt}\par\noindent
(1) 윗글의 문장 ``\textit{Any books in good condition are gladly accepted}''를, \underline{We를 주어로 하는 능동태} 문장으로 바꿔 쓰시오.
\par\vspace{4pt}\noindent
$\Rightarrow$ \rule{11cm}{0.4pt}
\vspace{10pt}\par\noindent
(2) 다음 문장을 우리말로 자연스럽게 해석하시오.
\par\vspace{2pt}\noindent
\hspace{1em}\textit{We would be grateful if you could join this meaningful event.}
\par\vspace{4pt}\noindent
$\Rightarrow$ \rule{\linewidth}{0.4pt}

\end{document}
